% frontier_mnop_framing_volume.tex
% Standalone compact-CY3 frontier: theorem-grade survivors and named
% obstruction objects against the formal-local mixed HT SFT on R^2 x C^2.

\documentclass[12pt]{amsart}
\usepackage[localtheorems]{raeez-math-template}
\newcommand\nc{\newcommand}
\newcommand\rc{\renewcommand}

\newcommand\nm{}

\newcommand\mc[1]{\mathcal{#1}}
\newcommand\mbb[1]{\mathbb{#1}}
\newcommand\mbf[1]{\mathbb{#1}}
\newcommand\mrm[1]{\mathrm{#1}}

\newcommand\image[1]{
  \begin{figure}[H]
    \centering
    \refstepcounter{figure}\label{fig:#1}
    \includegraphics[scale=.5]{#1}
  \end{figure}
}
\newcommand\imageopt[3]{
  \begin{figure}[H]
    \centering
    \ifstrequal{#2}{}
      {\includegraphics{#1}}
      {\includegraphics[scale=#2]{#1}}
    \ifstrequal{#3}{}{}{\caption{#3}}
    \label{fig:#1}
  \end{figure}}

\newcommand\mathsc[1]{\text{\normalfont\scshape#1}}
\providecommand{\Prim}{\operatorname{Prim}}
\providecommand{\Conn}{\operatorname{Conn}}

\newcommand\lb{\left[}  \newcommand\rb{\right]}
\newcommand\llb{\lb\lb}  \newcommand\rrb{\rb\rb}

\newcommand\lp{\left(}  \newcommand\rp{\right)}
\newcommand\llp{\lp\lp} \newcommand\rrp{\rp\rp}

\newcommand\twotuple[2]{\lp #1,#2\rp}
\newcommand\setpresent[2]{\{ #1 \mid #2\}}

\nc\Cat[1]{\mathbf{#1}}
\nc\G{G}
\nc\ringint{\mathcal{O}}
\nc\grk{k}
\nc\N{\mathbb{N}}
\nc\Z{\mathbb{Z}}
\nc\Q{\mathbb{Q}}
\nc\R{\mathbb{R}}
\nc\C{\mathbb{C}}
\nc{\FnS}{\grk\lp\Sigma\rp}
\nc\laurentpoly{\lb x,x^{-1}\rb}
\nc\laurentser{\llp~z\rrp}
\nc\fpowerser{\llb~z\rrb}
\nc\klocint{\grk\fpowerser}
\nc\Klocfrac{K\laurentser}
\nc\tensor{\otimes}
\nc\isom{\simeq}

\nc\ov[1]{/ {#1}}%'over' as in Scheme over k or galois extension over k

\nc\restr[2]{{
  \left.\kern-\nulldelimiterspace
  #1
  \vphantom{\big|}
  \right|_{#2}
  }}

\nc\shf[1]{\mathcal{#1}}
\nc\idshf[1]{\mathcal{#1}}

\nc\projline{\mathbb{P}^1}

\nc\K{\mathbb{K}}

\nc\polyring[2]{R}

\nc\classgrp[1]{C_{#1}}
\nc\F{\mathbb{F}}
\nc\idealgenby[1]{\langle#1\rangle}
\nc\D{\mathbf{d}}
\nc\pt{\mathbf{pt}}
\nc\Gr{\textrm{Gr}}

  \nc\Ga{\mathbb{G}_a}

  \nc\Gm{\mathbb{G}_m}

\nc\lpgrass{\mathcal{G}}
\nc\glpgrass{\mathcal{G}^P (I,Q)}

\nc\Flds{\Cat{Fields}}
\nc\cartdual{\mc{D}}

\nc\Mat[1]{\mathbf{Mat}_{#1\times#1}}
\nc\RH{\mathbb{\R}\textrm{H}}
\nc\tateshiftedby[1]{\left[#1\right]}
\nc\TateCat{\Cat{J_B}}
\nc\TateCatG{\Cat{J_G}}
\nc\B[1]{\mathbb{B} (#1)}
\nc\ptmod[1]{\mathbf{pt}\ov{#1}}
\nc\Ind{\mathbf{Ind}}
\nc\RHom{\mathrm{RHom}}
\nc\Hom{\mathrm{Hom}}
\nc\Rep{\mathrm{Rep}}
\nc\End{\mathrm{End}}
\nc\Map{\mathrm{Map}}
\nc\coH{\mathbf{H}}
\nc\hcoH[3]{\widehat{\coH^{#1}#2,#3}}
\nc\Ab{\Cat{Ab}}
\nc\Ker{\mathbf{Ker}}
\nc\Coker{\mathbf{Coker}}
\nc\dg{\textrm{DG}}
\rc\Vec{\Cat{Vec}}
\nc\VecF{\Cat{Vec}_{\grk}}
\nc\Coh{\Cat{Coh}}

\nc\Coind{\mathbf{Coind}}
\nc\Res{\mathbf{Res}}

\nc\Pic{\mathbf{Pic}}
\nc\AJ{\mathbf{AJ}}
\nc\divisor[1]{\mathbf{#1}}
\nc\divisoridshf[1]{\idshf{I_{\divisor{#1}}}}
\nc\IndProSch{\Cat{IndProSch}}
\nc\FinIndProSch{\Cat{FinIndProSch}}
\nc\Pro{\mathbf{Pro}}
\nc\IndPro{\Ind-\Pro}
\nc\FinIndSch{\Cat{FinIndProSch}}
\nc\A{\Cat{CommIndSch_{\grk}}}

\nc\Fun[2]{#1( #2 )}
\nc\Frac[1]{\mathbf{Frac}( #1 )}
\nc\curriedleftarg{\ldots}
\nc\Zp{\Z_p}

\nc\Qp{\mathbb{Q}_p}

\nc\Set{\Cat{Set}}

\nc\Hilb{\mathcal{H}\mathbb{ilb}}

\nc\disj{\sqcup}
\nc\bij{\leftrightarrow}
\nc\sur{\twoheadedrightarrow}
\nc\inj{\rightarrowtail}
\nc\injects{\inj}
\nc\fst{\ensuremath{2^{\text{\tiny{st}}}}}
\nc\snd{\ensuremath{3^{\text{\tiny{nd}}}}}
\nc\thrd{\ensuremath{3^{\text{\tiny{rd}}}}}
\nc\nth{\ensuremath{n^{\text{\tiny{th}}}}}
\nc\shfcoH[3]{\mathbf{H}^{#3}(#1,#2)}
\nc\hshfcoH[3]{\widehat{\shfcoH{#1,#2,#3}}}
\nc\Obj{\mathbf{Obj}}
\nc\grpring[2]{#1 \left[ #2 \right]}

\nc\Ext{\mathbf{Ext}}
\nc\Exts[3]{\Ext_{#1}^*\left(#2,#3 \right)}
\nc\dualmod[1]{#1^{*}}
\nc\LModCat[1]{#1-\Cat{LMod}}
\nc\RModCat[1]{#1-\Cat{RMod}}
\nc\indlim{\mathrm{lim}}

\nc\Schemes{\Cat{Sch}}

\nc\NoethSch{\Cat{NoethSch}}
\nc\NoethIntSch{\Cat{NoethIntSch}}

\nc\Sch[1]{\Cat{Sch}_{#1}}
\nc\IndSch[1]{\Cat{IndSch}_{#1}}

\nc\ModCat[1]{#1-\Cat{Mod}}

\nc\grkAlg{\Cat{\grk-Algebras}}
\nc\ZAlg{\Cat{\Z-Algebras}}

\nc\cc
{\ensuremath{\mathbf{CC}}}
\nc\Bun{\ensuremath{\mathrm{Bun}}}

\nc\Spec[1]{\mathrm{Spec} (#1)}
\nc\Specf[1]{\mathrm{Specf} (#1)}

\nc\mK{\mathbf{K}}
\nc\Bil{\mathbf{Bil}}
\nc\IndNSt{\Cat{Ind}-\textsc{N}-\Cat{Stacks}}

\nc\SL{\mathbf{SL}}
\nc\GL{\mathbf{GL}}
\nc\DMod{\Cat{DMod}}
\nc\IndSt{\Cat{IndStacks}}

\nc\Curve{\Sigma}
\nc\IxCHilb{\Hilb_{\Curve \times I}}
\nc\HilbCxI{\Hilb_{\Curve \times I}}
\nc\locLb{\mc{L}_{I,Q}}
\nc\Locvbpos{\{V_p\}_{p \in P}}

\nc\Tor{\mathbf{Tor}}
\nc\W{\mc{W}}
\nc\T{\mathbf{T}}
\nc\Hirz{\mbf{H}}
\nc\Tot{\mbf{Tot}}

\rc\P{\mathbb{P}}
\nc\Pbd{l_{\infty}}
\nc\blP{\widetilde{\mathbb{P}^2 \ov{\Z_2}}}
\nc\cM{\widetilde{\mathbb{M}(r,c)}}
\nc\M{\mathbb{M}(r,n)}
\nc\ef{\frac{\varepsilon_1}{\varepsilon_2}}
\nc\Vleft{\V{\sqrt{\frac{\ef}{\ef - 1}}}} \nc\Vright{\V{\sqrt{(1 - \ef)}}}
\nc\Verma[2]{\mathbf{V}_{#1,#2}}
\nc\VermaL[2]{\mathbf{L}_{#1,#2}}
\nc\V[1]{\mathbb{V}_{#1}}

\nc\lie[1]{\mathfrak{#1}}
\nc\su{\lie{su}}
\nc\gl{\lie{gl}}

\nc\g{\lie{g}}


% amsthm setup (preamble.tex provides these, but we're standalone here).
\newtheorem{thm}{Theorem}[section]
\newtheorem{prop}[thm]{Proposition}
\newtheorem{conj}[thm]{Conjecture}
\theoremstyle{remark}
\newtheorem{rmk}[thm]{Remark}
\renewcommand*{\theHthm}{\thesection.\arabic{thm}}
\renewcommand*{\theHprop}{\theHthm}
\renewcommand*{\theHconj}{\theHthm}
\renewcommand*{\theHrmk}{\theHthm}

% Local macros used only in this frontier file.
\providecommand{\E}{\mathrm{E}}
\providecommand{\PhiFA}{\Phi^{\mathrm{FA}}}
\providecommand{\DT}{\mathrm{DT}}
\providecommand{\PT}{\mathrm{PT}}
\providecommand{\GW}{\mathrm{GW}}
\providecommand{\Perf}{\mathrm{Perf}}
\providecommand{\Aut}{\mathrm{Aut}}
\providecommand{\End}{\mathrm{End}}
\providecommand{\cF}{\mathcal{F}}
\providecommand{\cO}{\mathcal{O}}
\renewcommand{\P}{\mathbb{P}}
\providecommand{\NAJ}{N_{\AJ}}
\providecommand{\srcdir}{\texttt{references/primary-sources/}}

\title[Compact CY$_3$ frontier]{Compact CY$_3$ frontier: theorem-grade
  survivors and named obstruction objects}

\begin{document}
\maketitle

\section*{Frontier surface}

The bound theorem of \texttt{main.tex} is local, on
\(\mathbb R^2_{\mathrm{top}}\times\mathbb C^2_{\mathrm{hol}}\): a
shifted cotangent Hamiltonian Lie algebra, its coordinate CE/PV
identity, principal-part currents, weighted Moyal/BV coefficients.  A
compact Calabi--Yau threefold \(X\) sits outside that surface.  The
descent obligation that would carry a local statement to compact \(X\) is
a matched-conventions theorem with explicit null-homotopies of the
period, anomaly, QME, convergence, hyperdescent, six-relation, and
operadic obstruction classes; no such theorem is supplied here, and
nothing below depends on the local Hamiltonian BF/Moyal calculation.

What survives at theorem grade on the compact CY\(_3\) frontier is
short.  The Maulik--Nekrasov--Okounkov--Pandharipande (MNOP) DT/PT
factorization is a scalar partition-function identity on a smooth
projective Calabi--Yau threefold.  On the quintic
\(Q\subset\mathbb P^4\), the Calabi--Yau volume form orients the normal
bundle of an \(S^3\)-Lagrangian and trivializes its rank-three
classifying class, leaving the integer-framing torsor and the
chain-level homotopy as a precisely named obstruction.  On
\(K3\times E\), Bridgeland--Maciocia plus Mukai--Poincar\'e produce a
derived equivalence to a Fourier--Mukai partner under a coprimality
hypothesis.  The Serre functor \(S_X=[3]\) and its pairing are an
instant categorical input on any smooth compact CY\(_3\).  Outside this
quartet, every claim names a precise missing object: the
\(\E_2\)-centre automorphism \(\sigma_X\), the chain-level framing
homotopy, the Hodge-normalisation datum \(\NAJ\), the global resurgent
sheaf for non-perturbative BCOV, the Serre-equivariant local descent
witnessing a quasi-NCCR.

Every external import here (CoHA, OSV/GV, Abel--Jacobi, Igusa/BKM, Vol
III chiral-volume, Kuznetsov HPD) is a quarantined comparison anchor in
the sense of \texttt{CLAUDE.md}.  Convention divergence with a sister
volume is load-bearing: it is reported, not silently reconciled.  The
source fixtures named throughout sit under \srcdir{} and are listed in
\S\ref{sec:source-anchors}; each fixture is an inventory of present and
missing source rows, not a constructed theorem.

\section{MNOP factorisation and the $\E_2$-centre obstruction}\label{sec:mnop}

Three enumerative theories on a smooth compact Calabi--Yau threefold
\(X\) ---
Donaldson--Thomas counts of ideal sheaves
\(Z_{\DT}(q)=\sum_n\DT_n(X)q^n\) (Thomas 2000; Behrend--Fantechi
2008),
Pandharipande--Thomas stable-pair counts
\(Z_{\PT}(q)=\sum_n\PT_n(X)q^n\) (Pandharipande--Thomas 2009), and
Gromov--Witten virtual-class generating series
\(Z_{\GW}(u)=\sum_g\GW_g(X)u^{2g-2}\) (Behrend--Fantechi 1997;
Li--Tian 1998) ---
produce the same complex numbers under the substitution \(-q=e^{iu}\)
and a Katz--Klemm--Vafa integrality wedge.  Why?  At the level of
scalar partition functions, Maulik--Nekrasov--Okounkov--Pandharipande
formulate this as a single MNOP package; the proven half is the DT/PT
factorisation by the MacMahon power off the degree-zero locus.

\begin{thm}[DT/PT factorisation on the quintic]\label{thm:mnop-quintic}
  Let \(Q\subset\P^4\) be a smooth quintic threefold (Euler
  characteristic \(\chi(Q)=-200\), Calabi--Yau by \(c_1(Q)=0\)).  The
  degree-zero DT generating function factors as the MacMahon power
  \(M(q)=\prod_{n\ge1}(1-q^n)^{-n}\) raised to \(\chi(Q)\), and the
  reduced DT series equals the PT series:
  \[
    Z_{\DT}(Q;q) \;=\; M(-q)^{\chi(Q)}\,Z_{\PT}(Q;q).
  \]
\end{thm}

\begin{proof}
  The factorisation \(Z_{\DT}=Z_{\DT,0}\cdot Z'_{\DT}\) with
  \(Z_{\DT,0}=M(-q)^{\chi(Q)}\) is Behrend--Fantechi's Hilbert-scheme
  formula applied to the degree-zero locus.  The reduced equality
  \(Z'_{\DT}=Z_{\PT}\) is Bridgeland's Hall-algebra theorem (with
  Toda's auxiliary wall-crossing route).  Smoothness of \(Q\) and
  \(\chi(Q)=-200\) close the arithmetic.
\end{proof}

\begin{rmk}[The Pandharipande--Thomas leg]
  Theorem~\ref{thm:mnop-quintic} is the proven scalar leg of the MNOP
  package on the quintic.  The PT/GW leg
  \(Z_{\PT}(Q;q)=Z_{\GW}(Q;u)|_{-q=e^{iu}}\) is conditional on KKV
  integrality for \(Q\); finite-genus quintic evidence requires the
  Huang--Klemm--Quackenbush bootstrap.  HKQ is a missing local source
  row (\S\ref{sec:source-anchors}).
\end{rmk}

The substitution \(-q=e^{iu}\) is not a tautology between two notations
for the same generating function.  At chain level, the DT, PT, and GW
sides come from three a priori distinct dualisable modules over a
factorisation-algebra image of \(\Perf(Q)\), each carrying its own
partition function; a single chain-level automorphism must respect all
three trace maps to produce the substitution.  The compact target datum
that names this object is
\[
  C_{\mathrm{tar}}(X)=
  \bigl(\cF_X,\; M_{\DT},M_{\PT},M_{\GW},\;
    \mathrm{Tr}_{\DT},\mathrm{Tr}_{\PT},\mathrm{Tr}_{\GW},\;
    \sigma_X\bigr).
\]
Here \(\cF_X=\PhiFA_3(\Perf(X))\) denotes the conjectural image of the
3-Calabi--Yau dg category of perfect complexes on \(X\) under a
factorisation-algebra functor of degree~\(3\); each \(M_\bullet\) is a
conjectural dualisable \(\cF_X\)-module attached to the named theory;
each \(\mathrm{Tr}_\bullet\) is a conjectural trace map recovering the
corresponding scalar generating function; and
\(\sigma_X\in\Aut_{\E_2}(Z_{\E_2}(\cF_X))\) is a conjectural
automorphism of the \(\E_2\)-centre
\(Z_{\E_2}(\cF_X)\) of \(\cF_X\) (the centre under
the Lurie \(\E_n\)-centre formalism, here \(n=2\) since
\(Z_{\E_2}\) controls braided-monoidal automorphisms) implementing the
MNOP/PT--GW substitution.  None of these objects is constructed in the
local mixed HT SFT theorem surface.

\begin{conj}[$\E_2$-centre MNOP package, conditional]
  \label{conj:e2-centre-mnop-package}
  Suppose \(C_{\mathrm{tar}}(Q)\) is constructed in conventions matched
  to the scalar fixture of Theorem~\ref{thm:mnop-quintic}.  Then there
  exists a unique
  \(\sigma_Q\in\Aut_{\E_2}(Z_{\E_2}(\cF_Q))\) effecting the substitution
  \(-q=e^{iu}\) on
  \((\mathrm{Tr}_{\DT},\mathrm{Tr}_{\PT},\mathrm{Tr}_{\GW})\), and
  rigidity holds: any two such automorphisms agree in the
  \(\E_2\)-centre cohomology of \(\cF_Q\).
\end{conj}

The construction problem decomposes into four named missing objects:
the factorisation-algebra functor \(\PhiFA_3\) on a 3CY dg category
(Joyce--Gross--Joyce--Tanaka Hall--vertex source vocabulary; Lurie centre
formalism); the dualisable module structure on
\(M_{\DT},M_{\PT},M_{\GW}\); the chain-level trace maps making each
\(Z_\bullet\) scalar; and a centrality homotopy
\(\sigma_X\colon\mathrm{Tr}_{\DT}\Rightarrow\mathrm{Tr}_{\PT}\) (and
analogously to \(\mathrm{Tr}_{\GW}\)) in the \(\E_2\)-centre.  The
rigidity statement is the vanishing of an obstruction class
\(\mathrm{Ob}_{\mathrm{UKD}}(C_{\mathrm{tar}})\) in the unique-Kuznetsov-descent
group of the centre, controlling \(\E_2\)-uniqueness of \(\sigma_X\)
under the prospective dg/derived descent.  Joyce--Upmeier and
Kinjo--Park--Safronov supply orientation and CoHA prerequisites; the
local mixed HT SFT does not.  Convention divergence with the Vol III
chiral-volume programme is recorded as a quarantined comparison gap and
is not silently reconciled.

\section{$S^3$-framing on the quintic}\label{sec:s3-framing}

Costello's chain-level open topological B-model with boundary on a
Lagrangian \(L\subset X\) requires a framing datum on \(L\): a chosen
trivialisation of the \(SO(\mathrm{rank}\, N_L)\)-torsor of orthonormal
frames of the normal bundle, compatible with the brane-line BV complex.
Per the discipline of \texttt{CLAUDE.md}, framing on \(S^3\) is not
canonical; every framing is a choice with stated dependence.  At a
Lagrangian \(L\cong S^3\hookrightarrow Q\) in the smooth quintic --- for
instance the real locus of a real deformation (Sheridan 2016;
Sheridan--Smith 2021) --- the topological framing obstruction lives in
\[
  \pi_0\bigl(\Map(S^3,SO(3))\bigr) \;=\; \pi_3(SO(3)) \;=\; \Z,
\]
the Hirzebruch \(\sigma\)-invariant torsor.  Two questions separate.
First: does the Calabi--Yau structure on \(Q\) orient \(N_{L/Q}\) and
make it topologically trivial, so that the framing torsor is in fact
\(\Z\)?  Second: which integer in that torsor, and which
BV/\(A_\infty\) chain homotopy on the brane complex, does Costello's
open B-model require?  The first is the proposition below; the second
is the named obstruction surface.

\begin{prop}[Normal-bundle orientation and topological triviality]
  \label{prop:s3-framing-quintic}
  Let \(L\cong S^3\subset Q\) be a Lagrangian in the smooth quintic
  threefold (Calabi--Yau by \(K_Q=\cO_Q\), with nowhere-vanishing
  holomorphic volume form \(\Omega_Q\in H^0(Q,K_Q)\)).  Then \(N_{L/Q}\)
  is canonically oriented and the oriented rank-three real bundle
  \(N_{L/Q}\to S^3\) is topologically trivial.  The orientation comes
  from
  \[
    \det T_L\otimes\det N_{L/Q} \;=\; \det(TQ|_L)
    \;\cong\; \cO_L
  \]
  via restriction of \(\Omega_Q\).
\end{prop}

\begin{proof}
  Since \(K_Q=\cO_Q\) and \(\Omega_Q\) is nowhere zero, its restriction
  is a nowhere-vanishing section of \(\det(TQ|_L)\).  The Lagrangian
  splitting \(TQ|_L=TL\oplus N_{L/Q}\) as real bundles gives
  \(\det(TQ|_L)_{\mathbb R}=\det(T_L)\otimes\det(N_{L/Q})\); parallelisability
  of \(S^3\) makes \(\det(T_L)\) canonically trivial, so
  \(\det(N_{L/Q})\) is trivial and \(N_{L/Q}\) is oriented.  Every
  oriented rank-three real bundle on \(S^3\) is trivial: the
  classifying map lies in \(BSO(3)\) and
  \(\pi_3(BSO(3))=\pi_2(SO(3))=0\).
\end{proof}

\begin{rmk}[The named missing chain-level datum]\label{rmk:framing-obstruction}
  Proposition~\ref{prop:s3-framing-quintic} closes the topological
  triviality of \(N_{L/Q}\), not the chain-level framing.  Two further
  data are required for the Costello open B-model on \((Q,L)\): an
  integer framing class \([f]\in\pi_3(SO(3))=\Z\) selecting a point in
  the trivialisation torsor (a choice, not canonical), and a
  Costello-compatible BV/\(A_\infty\) chain homotopy
  \(h_f\) trivialising the obstruction class
  \[
    \mathrm{Ob}_{\mathrm{frame}}(L,f)\in
    H^0\bigl(\Omega^\bullet(L)\otimes\End_{\mathrm{BV}}\bigr)
  \]
  to the equivariant lift of the brane propagator over the chosen
  framing.  The Calabi--Yau volume form supplies the topological
  orientation and triviality input; \(([f],h_f)\) is the named missing
  chain-level datum.
\end{rmk}

\section{Chiral-volume finite evidence on the quintic}\label{sec:chiral-volume}

If the genus-zero BPS counts \(n_N\) of degree-\(N\) rational curves on
a compact Calabi--Yau threefold grow at the rate predicted by the
period of a vanishing cycle, the sequence \(a_N=N^{-1}\log n_N\)
should converge.  To which value, and from which side?  The Vol III
chiral-volume conjecture proposes the Hodge-normalised Abel--Jacobi
period of a curve representing the class.  This is quarantined context
in the sense of \texttt{CLAUDE.md}: it lives outside the local mixed HT
SFT theorem surface and is not imported here as a consequence of any
local construction.  What can be tested locally is finite numerical
evidence on the quintic, the only compact CY\(_3\) where the BPS
sequence is currently known through \(N=10\).

Comparing local conditions:
\[
  \lim_{N\to\infty}\frac1N\log\bigl|Z_N(X,\beta)\bigr|
  \;\stackrel{?}{=}\;
  \tfrac1{2\pi}\,\bigl|\AJ(C_\beta)\bigr|
\]
in the quarantined Vol III statement, where \(Z_N=Z_N(X,\beta)\) is the
degree-\(N\) rational-curve count in class \(\beta\in H_2(X,\Z)\) and
\(\AJ(C_\beta)\in J^2(X)\) is the Abel--Jacobi period (the
intermediate Jacobian \(J^2(X)=H^3(X,\C)/(F^2H^3+H^3(X,\Z))\) carries the
Hodge-norm) of a representative curve \(C_\beta\).  Take \(X=Q\) the
smooth quintic, \(\beta=NH\) for \(H\in H_2(Q,\Z)\) the hyperplane class,
\(n_N=Z_N(Q,NH)\) the genus-zero BPS invariant.  The finite data
through \(N=10\) (BCOV bootstrap for \(N\le 9\) and the
Huang--Klemm--Quackenbush row for \(N=10\), conditional on the source
fixtures being supplied) is:
$$\renewcommand{\arraystretch}{1.1}
\begin{array}{r@{\;=\;}l}
  n_1 & 2875 \\
  n_2 & 609\,250 \\
  n_3 & 317\,206\,375 \\
  n_4 & 242\,467\,530\,000 \\
  n_5 & 229\,305\,888\,887\,625 \\
  n_6 & 248\,249\,742\,118\,022\,000 \\
  n_7 & 295\,091\,050\,570\,845\,659\,250 \\
  n_8 & 375\,632\,160\,937\,476\,603\,550\,000 \\
  n_9 & 503\,840\,510\,416\,985\,243\,645\,106\,250 \\
  n_{10} & 704\,288\,164\,978\,454\,686\,113\,488\,249\,750.
\end{array}$$
The sequence $a_N := N^{-1} \log n_N$ attains its minimum at $N=3$:
$$
a_1 \approx 7.964,\;
a_2 \approx 6.660,\;
a_3 \approx 6.525,\;
a_4 \approx 6.554,\;
a_5 \approx 6.613,
$$
$$
a_6 \approx 6.676,\;
a_7 \approx 6.733,\;
a_8 \approx 6.785,\;
a_9 \approx 6.832,\;
a_{10} \approx 6.873.
$$
The sequence decreases sharply through $N = 3$, then increases
monotonically: $a_{N+1} - a_N > 0$ for $N \in \{3, \dots, 9\}$, with
increments
\[
\{+0.029,\ +0.060,\ +0.062,\ +0.058,\ +0.052,\ +0.046,\ +0.041\}.
\]

\begin{prop}[Quintic chiral-volume finite evidence through $N=10$]
  \label{prop:chiral-volume-quintic}
  Take the displayed genus-zero quintic BPS table through \(N=10\) as
  input.  The sequence \(a_N=N^{-1}\log n_N\) attains its minimum at
  \(N=3\) and, on the tail \(N\in\{5,\dots,10\}\), satisfies the
  three-parameter law
  \[
    a_N \;=\; L \;+\; \alpha\,\frac{\log N}{N} \;+\; \frac{\beta}{N}
  \]
  with residual bounded by \(5\times10^{-6}\) and
  \[
    L \;\approx\; 7.5925,\qquad \alpha \;\approx\; -3.3204,
    \qquad \beta \;\approx\; 0.4477.
  \]
  The intercept \(L\) agrees, to three significant figures, with the
  raw mirror-period ratio at the conifold point \(z_c=5^{-5}\) of the
  Candelas--de la Ossa--Green--Parkes (CDGP) mirror quintic:
  \[
    |t(z_c)| \;=\;
    \biggl|\frac{\omega_1(z_c)}{\omega_0(z_c)}\biggr|
    \;\approx\; 7.5909,
  \]
  where the fundamental period
  \(\omega_0(z)=\sum_n\binom{5n}{n,n,n,n,n}z^n\) and the logarithmic
  period
  \(\omega_1(z)=\omega_0(z)\log z+5\sum_n\binom{5n}{n,n,n,n,n}(H_{5n}-H_n)z^n\)
  (with \(H_k\) the \(k\)-th harmonic number) solve the quintic
  Picard--Fuchs operator \(\mathcal L_{\mathrm{PF}}\) of
  \eqref{eq:picard-fuchs-cdgp} below.
\end{prop}

\begin{proof}
  Linear least squares on the six equations
  \(a_N=L+\alpha(\log N)/N+\beta/N\), \(N\in\{5,\dots,10\}\), gives the
  stated coefficients with maximum residual
  \(4.62\times10^{-6}<5\times10^{-6}\).  The fundamental and logarithmic
  periods solve
  \begin{equation}\label{eq:picard-fuchs-cdgp}
    \mathcal L_{\mathrm{PF}}\;=\;\theta^4-5z(5\theta+1)(5\theta+2)(5\theta+3)(5\theta+4),
    \qquad \theta=z\partial_z,
  \end{equation}
  with \(z=5^{-5}\psi^{-5}\) the natural mirror coordinate.  At
  \(z_c=5^{-5}\) the convergent power-series sum (verified against
  \texttt{scripts/quintic\_arithmetic\_rate\_checks.py}) yields
  \[
    \omega_0(z_c) \;\approx\; 1.0707258684,\qquad
    \omega_1(z_c) \;\approx\; -8.1277650532,
  \]
  giving \(|t(z_c)|\approx 7.5908925831\).
\end{proof}

\begin{rmk}[Three numerical paths]
  Three values are mutually compatible.  Direct tail extrapolation
  gives \(L\approx 7.59\); the raw mirror-period ratio gives
  \(|t(z_c)|\approx 7.5909\); the OSV attractor comparison value is
  \(5\log 5\approx 8.05\), a missing-source comparison row.  The fitted
  correction \(\alpha\approx -3.32<0\) forces approach to \(L\) from
  below, so the conjectural chiral-volume limit is not bounded above by
  \(\{a_1,a_2,a_3\}\); those low-degree points sit on the
  Schubert--Katz--Ellingsrud--Str\o mme Fano-geometric branch, below
  the expected large-\(N\) BPS regime.
\end{rmk}

\begin{rmk}[The named missing $\NAJ$ datum]
  \label{rmk:hodge-normalisation-aj}
  The quarantined Vol III conjecture
  \(\lim a_N=(2\pi)^{-1}|\AJ(C_H)|\) is stated against the Hodge metric
  on the intermediate Jacobian \(J^2(Q)\).  The raw CDGP mirror ratio
  \(t(z_c)=\omega_1(z_c)/\omega_0(z_c)\) is not an Abel--Jacobi norm:
  the periods \(\omega_0,\omega_1\) sit in a CDGP basis of solutions of
  \eqref{eq:picard-fuchs-cdgp} that has not been Hodge-normalised, and
  the \(2\pi i\) monodromy factor is not yet present.  The named
  missing object is a normalisation datum
  \[
    \NAJ\colon\;\bigl(\omega_0,\omega_1,\omega_2,\omega_3\bigr)
    \;\longmapsto\;\bigl(\varpi_0,\varpi_1,\varpi_2,\varpi_3\bigr)
  \]
  matching the CDGP period basis to the Hodge filtration on \(H^3\) and
  fixing the \(2\pi i\) monodromy convention.  Without \(\NAJ\), the
  finite agreement \(L\approx|t(z_c)|\) is a normalisation check, not a
  proof of the limit.  The OSV physical rate \(5\log 5\) is similarly a
  missing-source comparison value awaiting the OSV attractor row.
\end{rmk}

\section{Serre pairing as compact-CY categorical input}\label{sec:serre-nccr}

What is the categorical input that any compact CY\(_3\) descent of the
local theorem surface must respect, before any module, tilting, or
Hall-algebra structure is built?  The answer is the Serre functor and
its pairing.  Van den Bergh's noncommutative crepant resolution (NCCR)
is defined for singular Gorenstein CY\(_3\) via an Azumaya-like
generator of \(D^b\); it is Iyama--Wemyss in the toric case and a
consequence of the Orlov blow-up equivalence at \(K3\times E\).  On the
smooth quintic there is no crepant resolution to resolve, and the NCCR
question degenerates to the bare Serre datum.

\begin{prop}[Serre pairing on smooth compact CY$_3$]
  \label{prop:serre-nccr-substitute}
  Let \(X\) be a smooth compact Calabi--Yau threefold.  The Serre functor
  \(S_X\colon D^b(X)\to D^b(X)\) is the shift \(S_X=[3]\), and the
  Calabi--Yau pairing
  \[
    \RHom(F,G)\otimes\RHom(G,F)\;\longrightarrow\;\C[-3]
  \]
  is non-degenerate on \(\Perf(X)=D^b(X)\).
\end{prop}

\begin{proof}
  Smoothness gives \(\Perf(X)=D^b(X)\).  Serre duality is
  \(S_X=\otimes K_X[\dim X]\); the Calabi--Yau condition \(K_X=\cO_X\)
  reduces this to \(S_X=[3]\).  The pairing is the Serre trace of
  \(\RHom\otimes\RHom\to\RHom(F,F)\) followed by the Serre trace
  \(\RHom(F,F)\to\C[-3]\).
\end{proof}

\begin{rmk}[The named missing quasi-NCCR object]
  Proposition~\ref{prop:serre-nccr-substitute} fixes the Serre datum that
  a compact-CY descent must preserve.  The remaining global quasi-NCCR
  obstruction is the construction of a Serre-equivariant local
  algebraic or dg model
  \(\mathcal A_X\to\End_{D^b(X)}(\bigoplus_i F_i)\) generating
  \(D^b(X)\), with finite-dimensionality of \(\mathrm{Ext}^*(F_i,F_j)\)
  and a descent datum identifying the local model with a global
  endomorphism sheaf compatibly with \(S_X\).  This object exists in
  the singular Gorenstein cases above; it is not constructed for the
  smooth quintic by the present theorem surface.
\end{rmk}

\section{Fourier--Mukai partners for $K3 \times E$}\label{sec:fm-k3e}

A natural attempt to obtain a derived equivalence on a compact CY\(_3\)
is to take \(K3\times E\) and run Kuznetsov homological projective
duality.  This fails for a structural reason: HPD (Kuznetsov 2007)
takes as input a smooth projective variety with a very ample line
bundle inducing a non-trivial Lefschetz decomposition of \(D^b\), but
on any Calabi--Yau the Serre functor is a shift, and Bridgeland's
indecomposability theorem (Bridgeland 1999) forbids a non-trivial
semiorthogonal decomposition.  Both \(D^b(K3)\) and \(D^b(E)\) are
indecomposable, so the Lefschetz input on \(K3\times E\) is the
trivial rectangular structure \(\mathcal A_0=D^b(K3\times E)\), and
Kuznetsov Fano HPD does not apply.  What survives is a
Fourier--Mukai partner construction governed by Mukai-vector
coprimality, not a self-duality statement.

\begin{thm}[Fourier--Mukai equivalence in the isotropic \(K3\times E\) case]
  \label{thm:fm-self-duality-k3e}
  Let $S$ be a projective $K3$ surface with primitive polarisation
  $H_S$, and let $E$ be a principally polarised elliptic curve. Fix a
  Mukai vector $v = (r, \ell, s) \in \widetilde{H}(S, \Z)$ with $\ell
  \in \mathrm{NS}(S)$, Mukai pairing \(\langle v,v\rangle=0\), and
  $H_S$ chosen $v$-generic. Assume the coprimality condition
  \begin{equation}\label{eq:coprimality-star}
    \gcd(r,\, \ell \cdot H_S,\, s) = 1. \tag{$\star$}
  \end{equation}
  Let $M := M_{H_S}(v)$ be the moduli space of $H_S$-stable sheaves
  with Mukai vector $v$. Then:
  \begin{enumerate}
    \item $M$ is a smooth projective $K3$ surface (Mukai 1987;
      O'Grady 1997);
    \item there is an equivalence of $\C$-linear triangulated
      categories
      $$\Phi_{\mathcal{U}} :\ D^b(S) \xrightarrow{\,\sim\,} D^b(M),$$
      given by Fourier--Mukai transform along the universal sheaf
      $\mathcal{U}$ on $S \times M$ (Bridgeland--Maciocia 2001,
      \emph{Complex surfaces with equivalent derived categories},
      Thm.\ 1.1);
    \item combined with the Poincar\'e--Mukai autoequivalence
      $\mathrm{FM}_{\mathcal{P}} : D^b(E) \xrightarrow{\,\sim\,}
      D^b(\hat{E}) = D^b(E)$ (Mukai 1981), the product transform
      $$\Phi_{\mathcal{U}} \boxtimes \mathrm{FM}_{\mathcal{P}} :\
        D^b(S \times E) \xrightarrow{\,\sim\,}
        D^b(M \times E)$$
      is an equivalence of CY$_3$ categories preserving the Serre
      structure $S_{X} = [3]$.  When \(M\simeq S\) this is an
      untwisted Fourier--Mukai self-equivalence; in general it is a
      derived equivalence to the partner \(M\times E\).
  \end{enumerate}
\end{thm}

\begin{proof}
  (1) Under $(\star)$ and $v$-genericity of $H_S$, every
  $H_S$-semistable sheaf with Mukai vector $v$ is stable
  (no strictly semistables); the moduli $M$ is then a smooth projective
  holomorphic symplectic variety of dimension
  \(\langle v,v\rangle+2=2\).  Hence \(M\) is a \(K3\) surface.  For
  positive Mukai square the moduli space is higher-dimensional
  hyperk\"ahler and the product with \(E\) is no longer a CY$_3$ of the
  form asserted here.

  (2) The universal sheaf $\mathcal{U}$ exists on $S \times M$ under
  $(\star)$: $\gcd(r, \ell \cdot H_S, s) = 1$ kills the Brauer
  obstruction to a universal family (C\u{a}ld\u{a}raru 2000,
  Thm.\ 1.3.2). The Fourier--Mukai kernel
  $\mathcal{U} \in D^b(S \times M)$ satisfies the
  Bridgeland--Maciocia orthogonality criterion
  (Bridgeland--Maciocia 2001, \S 1) and induces an equivalence.
  Without $(\star)$, $\mathcal{U}$ exists only as a
  $p_M^*\alpha$-twisted sheaf for a Brauer class $\alpha \in
  \mathrm{Br}(M)$, yielding instead $D^b(S) \simeq D^b(M, \alpha)$
  (Huybrechts--Stellari 2005).

  (3) On $E$ with principal polarisation, $\hat{E} \cong E$
  canonically; the Poincar\'e bundle $\mathcal{P} \in D^b(E \times E)$
  is the kernel of $\mathrm{FM}_{\mathcal{P}}$ (Mukai 1981, Thm.\
  2.2). Tensoring kernels,
  $\mathcal{K} := \mathcal{U} \boxtimes \mathcal{P} \in D^b((S \times
  E) \times (M \times E))$ defines a Fourier--Mukai transform whose
  composition with its adjoint is the identity on each factor; the
  external-tensor convolution
  \((\Phi_1\boxtimes\Phi_2)\circ(\Psi_1\boxtimes\Psi_2)
  = (\Phi_1\circ\Psi_1)\boxtimes(\Phi_2\circ\Psi_2)\)
  for FM kernels (Huybrechts, \emph{Fourier--Mukai transforms in
  algebraic geometry}, Prop.\ 5.10) lifts the factor identities to the
  product.  The Serre functor is preserved by any exact equivalence of
  triangulated categories (uniqueness of the Serre functor up to natural
  isomorphism); since both \(\Phi_{\mathcal U}\) and
  \(\mathrm{FM}_{\mathcal P}\) are equivalences between CY categories
  of dimensions \(2\) and \(1\) respectively, the product Serre
  \(S_{S\times E}=S_S\boxtimes S_E=[2]\boxtimes[1]=[3]\) survives.
\end{proof}

\begin{prop}[Brauer-twisted dichotomy in the isotropic case]
  \label{prop:coprimality-dichotomy}
  Keep the isotropic hypothesis \(\langle v,v\rangle=0\) and drop
  $(\star)$ in Theorem~\ref{thm:fm-self-duality-k3e}. Then there
  exists a Brauer class $\alpha \in \mathrm{Br}(M)$, whose order
  divides $\gcd(r, \ell \cdot H_S, s)$, such that
  $$D^b(S \times E) \simeq D^b\bigl((M, \alpha) \times E\bigr),$$
  where the right-hand side denotes the derived category of
  $p_M^*\alpha$-twisted sheaves. The partner is
  \((M,\alpha)\times E\) with twisted \(K3\) factor; an untwisted
  kernel exists only under the coprimality hypothesis.
\end{prop}

\begin{proof}
  C\u{a}ld\u{a}raru 2000 Thm.\ 1.3.2 constructs the Brauer obstruction
  to the universal sheaf from the non-coprime Mukai vector; its order
  is the gcd. Huybrechts--Stellari 2005 (\emph{Equivalences of twisted
  K3 surfaces}) then identifies the derived equivalence as a twisted
  Fourier--Mukai. The elliptic factor contributes an untwisted
  $\mathrm{FM}_{\mathcal{P}}$.\qed
\end{proof}

\begin{rmk}[Why ``Kuznetsov Fano HPD'' is the wrong frame]
  The attempted ``Kuznetsov Fano HPD on $K3
  \times E$'' formulation requires a smooth Fano ambient
  with a non-trivial Lefschetz collection; $K3 \times E$ is CY$_3$,
  not Fano, and its $D^b$ is indecomposable. The correct replacement is
  the Fourier--Mukai partner construction
  (Theorem~\ref{thm:fm-self-duality-k3e}), in the isotropic
  Mukai-vector case and governed by coprimality $(\star)$, with the
  Brauer-twisted dichotomy
  (Proposition~\ref{prop:coprimality-dichotomy}) supplying the
  beyond-coprimality partner. The Kuznetsov--Perry 2019
  categorical-join / rectangular-HPD machinery does not upgrade
  Fourier--Mukai to Fano HPD here, precisely because neither
  $D^b(K3)$ nor $D^b(E)$ admits a non-trivial Lefschetz decomposition
  to feed into the rectangular construction.
\end{rmk}

\begin{rmk}[Vol III firewall: $\kappa$ values do not collide]
  Coprimality \((\star)\) is the primitive Mukai-vector hypothesis
  controlling the untwisted Fourier--Mukai kernel.  The quarantined Vol
  III chiral-bridge comparison records four \(\kappa\) values that
  must not be silently identified:
  \begin{itemize}
    \item \(\kappa_{\mathrm{cat}}(K3\times E)=\chi(\cO_{K3\times E})=0\),
      the total-space categorical Euler invariant by K\"unneth.  This is
      the relevant compact-CY\(_3\) total-space datum here.
    \item \(\kappa_{\mathrm{fibre}}(K3)=\chi(\cO_{K3})=2\), the fibre
      Euler datum.
    \item \(\mathrm{rk}\,\widetilde H^*(K3)=24\), the rank of the
      extended Mukai lattice
      \(\widetilde H^*(K3,\Z)=H^0\oplus H^2\oplus H^4\).
    \item \(\kappa_{\mathrm{ch}}^{\mathrm{Heis}}=3\) and
      \(\kappa_{\mathrm{BKM}}(\Phi_1)=5\) in the paramodular Igusa
      \(\Delta_5\) row, both quarantined Vol III / Igusa anchors with no
      matched-conventions descent into the present surface.
  \end{itemize}
  No equality among these values is asserted; the four values are
  recorded together as a firewall, not a comparison theorem.  Per
  \texttt{CLAUDE.md}, convention divergence with Vol III is load-bearing.
\end{rmk}

\section{Non-perturbative BCOV: Gopakumar--Vafa closure target}
\label{sec:bcov-np-closure}

The perturbative BCOV genus expansion on a compact Calabi--Yau threefold
is a divergent formal series of Gevrey class~\(1\): genus-\(g\)
amplitudes grow as \((2g-3)!/A^{2g-2}\) with \(A=\oint_{\gamma_c}\Omega\)
the period of a vanishing cycle at the nearest conifold singularity
(Shenker; Mari\~no et al.).  Borel summation on the physical ray fails;
a non-perturbative completion is required.  Two compatible candidates
present themselves --- a Gopakumar--Vafa BPS resummation, and a global
resurgent Stokes sheaf --- and the question for the compact-CY
frontier is whether these agree.

Costello--Li 2012 (arXiv:1201.4501) construct the classical BCOV field
theory and the perturbative BV quantisation problem, with elliptic-curve
uniqueness and Virasoro results within their stated scope.  No compact
CY\(_3\) all-genus theorem for the quintic amplitudes is asserted there
or here.  Write \(\{F_g^{\mathrm{BCOV}}(X,t)\}_{g\geq 0}\) for a
solution of the BCOV holomorphic anomaly equation with standard
large-radius and conifold-gap boundary data; the perturbative
generating function is
\[
  F^{\mathrm{pert}}(\lambda,t) \;=\; \sum_{g\geq 0}\lambda^{2g-2}F_g(X,t),
  \qquad \lambda=g_s.
\]
The Gopakumar--Vafa schema (Gopakumar--Vafa 1998) proposes the
non-perturbative form
\begin{equation}\label{eq:gv-expansion}
  F^{\mathrm{GV}}(\lambda,t) \;=\; \sum_{\beta>0}\sum_{g\geq 0}n_\beta^g
    \sum_{m\geq 1}\frac{1}{m}\bigl(2\sin(m\lambda/2)\bigr)^{2g-2}
    e^{-m\,\beta\cdot t},
\end{equation}
where \(n_\beta^g\in\Z\) are integer BPS invariants of the
class~\(\beta\).  Each fixed-\(\beta\) summand resums infinitely many
perturbative genera and produces Schwinger poles at \(\lambda=2\pi k/m\)
whose residues are pair-production amplitudes in the dual gauge
description.

\begin{conj}[GV non-perturbative closure target on the quintic]
  \label{conj:bcov-np-quintic}
  Assume the HKQ, CDGP, GV, OSV, and conifold-resurgence source rows
  recorded as missing in the quintic fixture are supplied and their
  normalisations match the scalar DT/PT fixture.
  Let $Q \subset \P^4$ be the smooth quintic. Let
  $F^{\mathrm{BCOV}}_g(Q, t)$ be the genus-$g$ solution of the BCOV
  holomorphic anomaly recursion with the CDGP large-radius input and the
  standard conifold gap boundary condition. Let $F^{\mathrm{GV}}$ be the
  Gopakumar--Vafa resummation~\eqref{eq:gv-expansion}, with BPS
  invariants $n_\beta^g(Q)$ computed by Huang--Klemm--Quackenbush 2009
  through genus $22$. Then in the Kahler cone $\Re t > t_c$, where
  $t_c$ is the image of the conifold locus $\psi = 1$ on the
  complex-structure moduli space of the mirror quintic $\widetilde Q$:
  \begin{enumerate}
    \item[(i)] \emph{(Conditional finite-genus evidence.)} For every
      $0 \leq g \leq 22$,
      \[
        F^{\mathrm{BCOV}}_g(Q, t) \;=\; [\lambda^{2g - 2}]\, F^{\mathrm{GV}}(\lambda, t)
      \]
      as formal power series in $e^{-\beta \cdot t}$; equality holds in
      $\Z[[e^{-t}]]$ after the Katz--Klemm--Vafa integrality normalisation.
      This item is not source-grade until the HKQ and PT/GW rows are
      locally present.
    \item[(ii)] \emph{(Convergence target.)}
      The genus-$0$ BPS series $\sum_\beta n_\beta^0 e^{-\beta t}$ is
      expected to have a nonempty convergence domain controlled by the
      conifold/entropy scale.  No coefficient-growth theorem proving this
      convergence is supplied here.
    \item[(iii)] \emph{(Conditional local conifold Stokes target.)} The trans-series
      completion of $F^{\mathrm{pert}}$ produced by the GV resummation
      is compatible near the conifold point $\psi = 1$ with the
      Candelas--de la Ossa--Green--Parkes 1991 mirror-period monodromy:
      the local Stokes normalization is expected to have \(S_c=1\), and
      the instanton action is \(A=2\pi i\,\tau_c\), where \(\tau_c\) is the
      vanishing-cycle period $\oint_{\gamma_c} \Omega_{\widetilde Q}$.
  \end{enumerate}
\end{conj}

\begin{rmk}[The named missing objects in (i)--(iii)]
  Each clause names a missing construction.  (i) would apply the
  Pandharipande--Thomas leg of MNOP (Theorem~\ref{thm:mnop-quintic}) to
  the BCOV genus-\(g\) amplitudes determined by the holomorphic anomaly
  recursion and boundary data; this is not a compact-CY\(_3\)
  Costello--Li all-genus theorem.  The CDGP prepotential, HKQ
  genus-\(22\) bootstrap, GV sine expansion, and PT/GW matching are
  the missing source rows.  (ii) is the convergence problem for the
  genus-zero BPS series: Castelnuovo-type genus bounds do not control
  the number \(n_\beta^0\) of rational curves, so finite-genus and
  finite-degree data alone cannot bound the coefficient growth of
  \(\sum_\beta n_\beta^0 e^{-\beta t}\); a coefficient-growth theorem is
  the named missing analytic object.  (iii) is the local conifold
  Stokes match: the leading factorial growth
  \(F_g^{\mathrm{pert}}(Q,t)\sim(2g-3)!/(2\pi i\tau_c)^{2g-2}\) along
  the Borel direction fixed by the vanishing cycle \(\gamma_c\) (with
  \(\tau_c=\oint_{\gamma_c}\Omega_{\widetilde Q}\)) would identify, by
  Strominger's conifold mechanism, the residue at this Borel
  singularity with a single hypermultiplet of mass \(|\tau_c|\); in
  \eqref{eq:gv-expansion} this is the leading-\(m\) contribution of a
  single BPS hypermultiplet with \(n_{\beta_c}^0=1\), \(\beta_c\in
  H_2(Q,\Z)\) the A-model class paired with \(\gamma_c\) under mirror
  symmetry, scoped to a neighbourhood of the conifold ray.  The
  \'Ecalle--Borel resurgent sum should give Stokes normalisation
  \(S_c=1\) across that ray.  The Strominger, Vafa/Schwinger,
  Shenker--Mari\~no, and Cecotti--Codesido--Mari\~no rows remain missing
  source obligations.
\end{rmk}

\begin{conj}[Conditional conifold trinity on the quintic]
  \label{conj:conifold-trinity-quintic}
  On the quintic $Q \subset \P^4$, the following three numerical data
  are expected to coincide as complex numbers (up to sign and overall
  $2\pi i$ normalisation), conditional on the missing conifold source
  rows and the cited resurgent normalization:
  \[
    A_{\mathrm{BCOV}} \;=\; A_{\mathrm{Schwinger}} \;=\; A_{\mathrm{period}},
  \]
  where
  \(A_{\mathrm{BCOV}}\) is the Shenker--Mari\~no instanton action of
  the perturbative BCOV series, read off as
  \(|F_g|/|F_{g+1}|\sim|A|^2/(2g-2)(2g-1)\);
  \(A_{\mathrm{Schwinger}}=2\pi i/m\) is the Schwinger-pole location of
  the Gopakumar--Vafa summand at multiplicity \(m\), rescaled by the
  class \(\beta_c\cdot t_c\); and
  \(A_{\mathrm{period}}=\oint_{\gamma_c}\Omega_{\widetilde Q}\) is the
  vanishing-cycle period at the conifold locus \(\psi=1\) on the mirror
  quintic.  Equality of all three forces the conifold Stokes
  normalisation \(S_c=1\) and a single BPS hypermultiplet
  \(n_{\beta_c}^0=1\).
\end{conj}

\begin{rmk}[Source obligations for the conifold trinity]
  Equality \(A_{\mathrm{BCOV}}=A_{\mathrm{period}}\) requires a
  Huang--Klemm--Quackenbush/Shenker--Mari\~no large-order row.  Equality
  \(A_{\mathrm{Schwinger}}=A_{\mathrm{period}}\) requires the
  Vafa/Schwinger pair-production and Strominger conifold rows.  The
  resurgent normalisation \(S_c=1\) requires the
  Cecotti--Codesido--Mari\~no row.  No global Stokes-sheaf construction
  is claimed.
\end{rmk}

\begin{rmk}[The named missing global resurgent sheaf]
  Conjecture~\ref{conj:bcov-np-quintic}(iii) Borel-sums along the
  physical ray and matches at the single conifold Stokes wall.  The full
  non-perturbative completion requires matching at every Stokes wall of
  the quintic moduli space: large complex structure (where
  \eqref{eq:gv-expansion} is the target definition), the conifold (where
  Strominger's hypermultiplet pairs with the vanishing cycle), and
  Landau--Ginzburg / Gepner (where the HKQ bootstrap supplies finite-genus
  data).  The named missing object is a global resurgent sheaf
  \(\mathcal R^{\mathrm{np}}\) on the moduli space whose unique
  trans-series \(F^{\mathrm{np}}(\lambda,t)\in\Gamma(\mathcal M,\mathcal R^{\mathrm{np}})\)
  has Stokes jumps matching the BCOV monodromy representation on
  \(H^3(\widetilde Q,\Z)\).  This sheaf is not constructed in the
  present surface, neither globally on the quintic moduli space nor in
  any compact CY\(_3\) generalisation.
\end{rmk}

\begin{conj}[Mirror-period recovery from non-perturbative BCOV]
  \label{conj:mirror-period-recovery}
  Assume the missing CDGP and GV source rows are supplied in conventions
  matched to those of Proposition~\ref{prop:chiral-volume-quintic}, with
  the mirror coordinate \(z=5^{-5}\psi^{-5}\).  Let
  \(\varpi_0(z),\varpi_1(z),\varpi_2(z),\varpi_3(z)\) be the
  Hodge-normalised CDGP 1991 fundamental periods of the mirror quintic
  \(\widetilde Q\), solutions of the Picard--Fuchs operator
  \(\mathcal L_{\mathrm{PF}}\) of \eqref{eq:picard-fuchs-cdgp}.  The
  genus-zero prepotential extracted from
  \(F^{\mathrm{GV}}(\lambda,t)=\lambda^{-2}\mathcal F_0(t)+O(\lambda^0)\)
  with mirror map \(t=\varpi_1/\varpi_0\) should satisfy
  \[
    \mathcal F_0(t) \;=\;
    \frac{1}{\varpi_0^{\,2}}\bigl(\varpi_0\varpi_3-\varpi_1\varpi_2\bigr)
    + \text{(classical polynomial)},
  \]
  with no instanton corrections beyond those encoded in
  \(\{n_\beta^0\}_\beta\).  The first three BPS values
  \(n_1^0=2875,\;n_2^0=609\,250,\;n_3^0=317\,206\,375\) come from the
  local BCOV table; the CDGP period/prepotential row and the
  Schubert--Katz--Ellingsrud--Str\o mme historical mirrors are missing
  source obligations.
\end{conj}

\begin{rmk}[Relation to the chiral-volume statement]
  Conjecture~\ref{conj:mirror-period-recovery} is the genus-zero
  partner of the chiral-volume statement
  (Proposition~\ref{prop:chiral-volume-quintic}) under the named missing
  \(\NAJ\) normalisation: with \(\NAJ\) supplied, the Hodge-normalised
  Abel--Jacobi period coincides with the mirror-coordinate modulus
  \(|t(z_c)|\), and the conjectural chiral-volume limit \(L\) becomes
  the logarithmic period
  \(|\varpi_1(z_c)/\varpi_0(z_c)|\) of the CDGP family at the conifold
  attractor.  Finite evidence is the HKQ genus-\(22\) data once its
  source row is supplied; higher-genus extension requires the resurgent
  completion of Conjecture~\ref{conj:bcov-np-quintic}(iii).
\end{rmk}

\section{Theorem-grade survivors and the named obstruction objects}
\label{sec:summary}

Four statements survive at theorem grade on the compact-CY\(_3\)
frontier in the present convention layer.  The DT/PT MacMahon
factorisation
\(Z_{\DT}(Q;q)=M(-q)^{\chi(Q)}Z_{\PT}(Q;q)\) on the smooth quintic
(Theorem~\ref{thm:mnop-quintic}) is a scalar partition-function
identity.  The normal-bundle orientation and topological triviality of
\(N_{L/Q}\) for an \(S^3\)-Lagrangian
(Proposition~\ref{prop:s3-framing-quintic}) reduces the framing torsor
to \(\Z\).  The Serre functor and pairing on smooth compact CY\(_3\)
(Proposition~\ref{prop:serre-nccr-substitute}) provide the categorical
Calabi--Yau datum \(S_X=[3]\).  The Fourier--Mukai equivalence on
\(K3\times E\) under primitive Mukai-vector coprimality
(Theorem~\ref{thm:fm-self-duality-k3e}) produces a derived equivalence
preserving \(S_X=[3]\); the Brauer-twisted dichotomy
(Proposition~\ref{prop:coprimality-dichotomy}) handles the non-coprime
case.

The conditional finite numerical evidence on the quintic
(Proposition~\ref{prop:chiral-volume-quintic}) sharpens this picture by
exhibiting the three-way numerical compatibility
\[
  L \;\approx\; 7.5925 \;\approx\; |t(z_c)| \;\approx\; 7.5909
\]
between the tail-fit intercept of \(a_N=N^{-1}\log n_N\) on the quintic
hyperplane class through \(N=10\) and the raw CDGP mirror-period ratio
at the conifold point \(z_c=5^{-5}\), against the missing OSV rate
\(5\log 5\approx 8.0472\) which it does not yet match.

What forces the next step is the precise list of named missing objects.
The MNOP package needs a chain-level
\(\sigma_Q\in\Aut_{\E_2}(Z_{\E_2}(\cF_Q))\) implementing the
substitution \(-q=e^{iu}\), with rigidity expressed by the vanishing of
\(\mathrm{Ob}_{\mathrm{UKD}}(C_{\mathrm{tar}})\).  The chain-level
\(S^3\)-framing needs an integer class \([f]\in\pi_3(SO(3))\) and a
Costello-compatible BV chain homotopy \(h_f\); per
\texttt{CLAUDE.md}, framing is a choice, not canonical.  The
chiral-volume statement needs the Hodge-normalisation datum
\(\NAJ\colon(\omega_0,\omega_1,\omega_2,\omega_3)\mapsto(\varpi_0,\varpi_1,\varpi_2,\varpi_3)\)
matching the CDGP basis to the Hodge filtration on \(H^3\).  The
quasi-NCCR descent needs a Serre-equivariant local algebraic or dg
model generating \(D^b(X)\).  The non-perturbative BCOV completion
needs a global resurgent sheaf \(\mathcal R^{\mathrm{np}}\) on the
moduli space whose Stokes jumps match the BCOV monodromy
representation on \(H^3(\widetilde Q,\Z)\).

Each of \(\sigma_Q\), \(([f],h_f)\), \(\NAJ\), the quasi-NCCR datum,
and \(\mathcal R^{\mathrm{np}}\) is a precise construction problem,
not a vague target.  None is supplied by the local mixed HT SFT
theorem surface of \texttt{main.tex}; convention divergence with the
Vol III chiral-volume programme, the Igusa \(\Delta_5\) bridge, and the
OSV attractor rate is recorded as a quarantined comparison gap, not
silently reconciled.

\section{Source anchors}\label{sec:source-anchors}

The fixtures referenced throughout are markdown rows under
\srcdir; each is an inventory of present and missing source rows, not
a constructed theorem.  In one place:
\begin{itemize}
  \item \texttt{mnop-pt-dt-source-anchors-2026-04-30.md} ---
    Donaldson--Thomas, Pandharipande--Thomas, Maulik--Nekrasov--Okounkov--Pandharipande,
    Bridgeland Hall-algebra, Behrend--Fantechi degree-zero MacMahon,
    Toda wall-crossing.  Anchors Theorem~\ref{thm:mnop-quintic}.
  \item \texttt{quintic-bcov-gv-source-anchors-2026-04-30.md} ---
    BCOV/Costello--Li framework, BCOV degree \(1,\dots,9\) quintic
    table, conifold vocabulary.  Records HKQ \(N=10\), CDGP, GV,
    OSV, Strominger, Vafa/Schwinger, Shenker--Mari\~no, and
    Cecotti--Codesido--Mari\~no rows as missing.
    Anchors Proposition~\ref{prop:chiral-volume-quintic} and
    Conjectures~\ref{conj:bcov-np-quintic},
    \ref{conj:conifold-trinity-quintic},
    \ref{conj:mirror-period-recovery}.
  \item \texttt{coha-center-source-anchors-2026-04-30.md} --- Lurie
    centre formalism, Joyce/Gross--Joyce--Tanaka Hall--vertex
    structures, Joyce--Upmeier and Kinjo--Park--Safronov
    orientation/CoHA prerequisites.  Source vocabulary for
    Conjecture~\ref{conj:e2-centre-mnop-package}; no constructed
    \(\cF_X\), trace maps, \(\sigma_Q\), \(\E_2\)-rigidity, or
    \(\mathrm{Ob}_{\mathrm{UKD}}\)-null-homotopy.
  \item \texttt{compact-cy-source-fixture-control-surface-2026-04-30.md}
    --- non-CoHA compact-CY domain fixture routing for the remaining
    historical and Hodge-theoretic source rows.
\end{itemize}

\end{document}
